\begin{center}
\begin{tikzpicture}\begin{axis}[width=.77\linewidth,height=8.2cm,xmin=-.5,xmax=19.5,ymin=-.5,ymax=19.5,y dir=reverse,xtick={0,1,2,3,4,5,6,7,8,9,10,11,12,13,14,15,16,17,18,19},ytick={0,1,2,3,4,5,6,7,8,9,10,11,12,13,14,15,16,17,18,19},xticklabels={3,16,0,4,10,6,7,2,18,5,8,15,19,13,9,14,12,11,1,17},yticklabels={alt.atheism,comp.graphics,comp.os.ms-windows.misc,comp.sys.ibm.pc.hardware,comp.sys.mac.hardware,comp.windows.x,misc.forsale,rec.autos,rec.motorcycles,rec.sport.baseball,rec.sport.hockey,sci.crypt,sci.electronics,sci.med,sci.space,soc.religion.christian,talk.politics.guns,talk.politics.mideast,talk.politics.misc,talk.religion.misc},tick label style={font=\scriptsize},xlabel={Кластер $k$-means},colorbar,point meta min=0,point meta max=8,colormap={density}{color=(white) color=(pine)},colorbar style={ytick={0,2,4,6,8}}]
\addplot[only marks,scatter,scatter src=explicit,mark=square*,mark size=3.6pt] table[x=x,y=y,meta=n]{ng-heatmap.dat};
\end{axis}\end{tikzpicture}
\par\small Число тестовых писем каждой рубрики в каждом кластере. В строке восемь писем. Для чтения рисунка столбцы переставлены по максимальному диагональному совпадению. Перестановка используется только при визуализации.
\end{center}
Два наиболее многочисленных тестовых кластера иллюстрируют связь названия с рубрикацией: кластер 14 «Christian Faith Persistence» (группа с наибольшим числом писем --- \texttt{soc.religion.christian}, 7 из 17) и кластер 10 «Controller Bus Speed Mismatch» (группа с наибольшим числом писем --- \texttt{comp.sys.ibm.pc.hardware}, 5 из 14).
